% Section 8 -- Discussion
\section{Discussion}
\label{sec:discussion}

\subsection{The direction result is a coverage result}

A leak is where a directional claim is most tempting and least
supportable. A leak presents as a path longer than the valley-free
expectation, a control-plane property, and its traffic consequence is a
redistribution among paths that often leaves the aggregate unchanged. The
evidence that identifies a leak---AS-relationship distance, the affected
prefix---is therefore control-plane evidence, and the cross-domain test
that would confirm it is a data-plane test that the traffic arm cannot
supply when the aggregate does not move. This is the regime V3 exists to
police: if the two series do not share the affected prefix, a timing
coincidence between them is not a causal claim, and
\texttt{provenance\_disjoint} is the honest verdict. The practical
consequence is a division of labour: Phase~1 classifies leaks, because it
reads path structure directly, and Phase~2 is an instrument for the
minority of events where a data-plane effect is measurable.

\subsection{The division of labour with language models}

The audit in \S\ref{sec:eval-audit} is the empirical basis for a design
rule we would otherwise assert on principle. A language model that
recovers 83\,\% of ASNs and 100\,\% of verdicts but 0\,\% of traffic
percentages is a competent \emph{renderer} of a decision and an
incompetent \emph{maker} of one. The architecture enforces the division
structurally---the model receives the verdict and the evidence table, and
its output is not read back into the decision---rather than by prompting,
because a soft constraint in the decision path is a reliability risk that
a model upgrade can convert into a regression. The cost is that a
narrative is only as informative as the evidence table, and on the
windows that stop at the first veto that table is mostly empty; a report
that names the veto that failed is less satisfying than a confident
paragraph, and it is the report the evidence supports.

\subsection{Threats to validity and what would change our mind}

\textbf{Construct.} The ground-truth labels are curated from published
accounts, and curation favours incidents that were reported. This
inflates the share of high-severity, clearly-caused events, which is
consistent with the gap between the curated set (76.5\,\%) and the
held-out set (61.5\,\%).

\textbf{Internal.} The label-leak audit removed \texttt{event\_type}
from three detectors and changed their outputs; every number in this
paper is computed after that audit. A second consideration is path
heterogeneity: the replay path that produces the Phase-2 numbers reads a
cache's full event streams where they exist and falls back to the merged
timeline where they do not, so some evaluated windows carry richer
evidence than the live path would supply at the same instant. That
discrepancy favours coverage: a fuller stream gives the estimator more
non-empty buckets to test than the live path supplies, which is the
opposite direction from the abstention result, and \S\ref{sec:eval-shape}
quantifies it on a window where both shapes are available.

\textbf{External.} Thirty curated events and seventeen agent-corpus
events are small evaluations, and the selective accuracy of the
directional layer rests on very few
commitments. We report counts rather than a rate so that a reader can
apply their own prior over how much two events should move the estimate.
The leak class is a third limitation: the corpus cannot separate leaks
from outages on physical grounds, so the leak row of
Table~\ref{tab:agent-gt} is a statement about the collection rather than
about the method. The agent-level comparison is also under-powered: at
$n=17$ the $13$ against $11$ gap between the agent pipeline and the rules
arm is suggestive rather than significant on its own, and we do not claim
it is; the per-class movement it produces is the defensible claim.

Three results would overturn the abstention-heavy design. A channel that
carried $R\to T$ signal on substantially more of the windows would
justify a lower floor, because the constraint would then be the evidence
rather than the arithmetic. An evaluation on events with better traffic
coverage---the held-out set's five uncovered ASes are the cautionary
example---would test whether the $R\to T$ gap is a sampling artefact or a
property of the data. And a leak evaluation with cross-domain
confirmation, from a third vantage point or an operator's own traffic
split, would test whether the leak blind spot is intrinsic to aggregate
traffic or specific to our single-series measurement. None of these is a
small experiment, which is why we report the state we can defend rather
than the state we would prefer.

\subsection{Ethical considerations}
\label{sec:ethics}

This study uses public data only: BGP updates from RIPE RIS, traffic
aggregates from Cloudflare Radar, and incident reports that network
operators published themselves. It involves no human subjects, no
end-user records, and no private operator telemetry; the traffic series
are aggregates over prefix sets, and the case study reproduces only what
the public report and the recorded artefact already contain
(\S\ref{sec:case}). The ethical risk we weighed is therefore not privacy
exposure but misattribution: a system that names the wrong AS can
misdirect an operator's response. That risk is why every abstention
carries a machine-readable reason rather than a low-confidence cause
(\S\ref{sec:method-gate}), and why \S\ref{sec:eval-errors} reports each
error individually rather than as an accuracy figure alone.
